%%%%%
%
\section{Intensit� de pr�cipitation pour un pluviom�tre Hasse}
\itemstyle{Op\'erateur~:} {\sf Rain\_rate\_Hasse}
\index{Rain\_rate\_Hasse}

%%%%%OBJET%%%%
\itemstyle{Objet~:}
%Algorithme de calcul d'une intensit� de pr�cipitation sous influence 
%d'un vent : Les param�tres d'entr�e sont les valeurs de 3 compteurs 
%(sup�rieur, lat�ral et an�mom�tre) � la m�me fr�quence que le vecteur
%de sortie.

Calculer une intensit� de pr�cipitation � partir des quantit�s d'eau
r�colt�es sur deux collecteurs (sup�rieur et lat�ral), en fonction
de la vitesse du vent mesur�e pr�s du pluviom�tre.


%%%%%ARGUMENTS%%%%
\bigskip
%%%%%ENTREES%%%%
\begin{tabarg}
\multicolumn{3}{l}{\titretab {Entr\'ees}}\\
\hline
$U$ & Vecteur, \texttt{float} & Vitesse du vent mesur�e par l'an�mom�tre associ� (m/s) \\
$Cs$ & Vecteur, \texttt{float} & R�sultat du comptage sup�rieur \\
$Cl$ & Vecteur, \texttt{float} & R�sultat du comptage lat�ral \\
unit & Vecteur, \texttt{string} & Unit� pour le param�tre de sortie \\
\end{tabarg}
\medskip

%%%%%SORTIES%%%%
\begin{tabarg}
\multicolumn{3}{l}{\titretab {Sorties}}\\
\hline
$rr$ & Vecteur, \texttt{float} & Intensit� de pr�cipitation\\
\end{tabarg}




%%%%%FORMULE%%%%
\itemstyle{Formule :}

%%Rotat_uv
%
%Algorithme utilis� :
%
%En entr�e :
%U	m/s	 Vitesse du vent mesur�e par l'an�mom�tre associ�
%Cs    	   	 R�sultat du comptage sup�rieur
%Cl   	 	 R�sultat du comptage lat�ral
%En sortie
%rr    	Mm/h	 Intensit� de pr�cipitation
%
Valeurs des coefficients utilis�s :
\begin{center}
% use packages: array
\begin{tabular}{l l}
a1 = 2.9562289E-02 & a2 = 9.8437145E-03 \\ 
a3 = -3.9051531E-03 & a4 = -4.8168935E-03
\end{tabular}
\end{center}

Calcul de l'intensit� de pr�cipitation pour la surface sup�rieure de collection :

%%\usepackage{amsmath}
\begin{equation*}
\begin{aligned}
rr_{sup} = &
\begin{cases}
0 & \text{si $Cs = 0$,} \\
krr_{sup} . Cs & \text{si $Cs > 0$.}\\
\end{cases} \\
& \text{avec } krr_{sup} =  exp ( a1 + a2 . U + a3 . log (Cs) + a4 . log (Cs) . U )
\end{aligned}
\end{equation*}

\bigskip
Si $U$ n'est pas nul, calcul de l'intensit� de pr�cipitation pour la surface lat�rale de collection :
\begin{equation*}
\begin{aligned}
rr_{lat} = &
\begin{cases}
0 & \text{si $Cl = 0$,} \\
20.9032 . lwc^{1.1264}
&  \text{si $Cl > 0$.}
\end{cases} \\
& \text{avec } lwc = \frac{Cl}{U} . ( 1 + 0.0015 . U^{2})
\end{aligned}
\end{equation*}
%\begin{eqnarray*}
%			klwc =  1 + 0.0015 . U^{2}\\
%%			lwc = \frac{Cl}{U} . klwc\\
%			rr_{lat} =  20.9032 . lwc^{1.1264})\\
%\end{eqnarray*}

\bigskip
Calcul de l'intensit� de pr�cipitation totale en fonction du vent :
\begin{equation*}
rr =
\begin{cases}
rr_{sup} & \text{si }  U < 9 m/s \\
\left( 1 - \frac{U - 9.}{11. - 9.}\right) . rr_{sup} + \left(\frac{U - 9.}{11. - 9.}\right) . rr_{lat}  & \text{si } U \leq 11 \text{ et } U \geq 9 m/s \\
rr_{lat} & \text{si }  U > 11 m/s \\
\end{cases} 
\end{equation*}



%%%%%COMMENTAIRE%%%%
\itemstyle{Commentaire~:}\\
Un premier traitement des donn�es brutes (comptage) a permis d'obtenir des taux d'accroissement :
Cs= 0.2972*(Vi-Vi-1)*120/10\\
Cl= 0.1607*(Vi-Vi-1)*120/10\\
U= (Vi-Vi-1)*0.1/10\\
On a moyenn� sur 10 minute pour "lisser" et �viter les retours � " z�ro sur 10 secondes, puis on a appliqu� l'algorithme\\
Remarque : On a proc�d� � plusieurs essais pour la moyenne, (sur 1, 2, 5 puis 10 minutes) car il n'existe pas d'indications du constructeur sur le fonctionnement de l'appareil. Les meilleurs r�sultats, comparativement � l'ORG sont obtenus avec 10 minutes.\\
\\
\underline{Comparaison � un autre pluviom�tre optique de type ORG :}\\
\textbf{Attention}, du fait de la comptabilisation des gouttes sur la surface lat�rale, les intensit�s Hasse sont sup�rieures � celles de l'ORG par vent cons�quent.

%%%%%AUTEUR%%%%
\itemstyle{Auteur~:}
CNRM/GMEI/TRAMM

%%%%%REFERENCE%%%%
\itemstyle{R\'ef\'erences~:}

Manuel du constructeur (Eigenbrodt) \cite{Hasse}.

